\documentclass[11pt,a4paper]{article}

\usepackage[utf8]{inputenc}
\usepackage[T1]{fontenc}
\usepackage{lmodern}
\usepackage[english]{babel}
\usepackage{geometry}
\usepackage{array}
\usepackage{longtable}
\usepackage{booktabs}
\usepackage{tabularx}
\usepackage{ragged2e}
\usepackage{enumitem}
\usepackage{xcolor}
\usepackage{hyperref}
\usepackage{titlesec}
\usepackage{fancyhdr}
\usepackage{listings}
\usepackage{graphicx}
\usepackage{float}

\geometry{margin=2.3cm}
\setlength{\parskip}{0.5em}
\setlength{\parindent}{0pt}
\renewcommand{\arraystretch}{1.2}

\pagestyle{fancy}
\fancyhf{}
\lhead{INTISAT -- ISM Payload}
\rhead{Rev. 0.1 Draft}
\cfoot{\thepage}

\newcolumntype{Y}{>{\RaggedRight\arraybackslash}p{11cm}}
\newcolumntype{L}[1]{>{\RaggedRight\arraybackslash}p{#1}}

\title{\textbf{INTISAT -- ISM PAYLOAD}\\
\large Complete Payload Architecture and Mission Engineering Document}
\author{Document ID: INTISAT-ISM-PAY-001}
\date{Revision 0.1 Draft --- 2026-04-09}

\begin{document}

\maketitle
\thispagestyle{fancy}

\begin{center}
\textbf{Classification:} Engineering Working Document
\end{center}

\section{Executive Summary}

The INTISAT Image Sensor Microscopy (ISM) payload is a CubeSat-class in-situ computational microscope implementing Fourier Ptychographic Microscopy (FPM), a computational imaging technique that reconstructs a high-resolution, complex-valued image (amplitude plus phase) from a sequence of lower-resolution intensity measurements taken under angularly diverse illumination.

The concept is validated at bench level. The software pipeline (capture $\rightarrow$ FPM reconstruction $\rightarrow$ super-resolution via RealESRGAN) runs on a Raspberry Pi 4 and produces demonstrably enhanced images. The repository contains approximately 18~GB of real experimental data, production-grade Python code, and extensive sensor/algorithm documentation. No hardware design files (PCB, schematic, or Gerber) exist in the repository.

This document closes the gap between the bench-level demonstrator and a flight-credible CubeSat payload by:
\begin{itemize}[leftmargin=1.5em]
    \item defining the hardware architecture and board partitioning,
    \item providing mission-level parameter tables and budgets,
    \item proposing an onboard versus ground processing split,
    \item identifying the differentiators with respect to prior work, and
    \item specifying the validation and test plan required for credibility.
\end{itemize}

\textbf{Recommended architecture:} Raspberry Pi Compute Module 4 (CM4) on a custom carrier board, with a separate sensor board and a microcontroller supervisor (Option D).

\section{Repository Audit}

\subsection{Audit Summary}
\begin{longtable}{L{4cm} L{3cm} L{7cm}}
\toprule
\textbf{Category} & \textbf{Status} & \textbf{Evidence} \\
\midrule
Capture code & Production-ready & \texttt{capturaclaude.py}, \texttt{captura\_circulo\_movil\_tftcolorSBARRIDO.py} \\
FPM reconstruction & Production-ready & \texttt{fpm\_procesador\_tft.py}, \texttt{fpm\_procesador\_tft\_fixed.py}, \texttt{fpm-py/} library \\
Super-resolution (ESRGAN) & Implemented / tested & \texttt{run\_esrgan\_fpm.py}, \texttt{rrdbnet\_standalone.py} \\
Metadata system & Implemented & JSON schema in capture scripts \\
Sensor documentation & Comprehensive & \texttt{CMOS\_SENSOR\_OV5647\_TECHNICAL\_GUIDE.md} \\
Algorithm documentation & Comprehensive & \texttt{TECHNICAL\_REPORT\_FPM\_WORKFLOW.md} \\
Real experimental data & Present ($\sim$18 GB) & \texttt{data/raw/}, \texttt{data/processed/}, \texttt{data/calibration/} \\
Hardware design (PCB/SCH) & Missing & Directory stubs only, no KiCad/Eagle/Gerber \\
Bus subsystems & Stubs only & ADCS, EPS, OBC, TCS all \texttt{.gitkeep} \\
Ground segment software & Stub & \texttt{ground/mcs/}, \texttt{ground/station/} empty \\
Link budget simulation & Stub & \texttt{sim/link\_budget/} empty \\
Orbital dynamics sim & Stub & \texttt{sim/dynamics/} empty \\
Power budget & Missing & No power analysis found \\
Thermal analysis & Missing & No thermal model found \\
Calibration data & Sparse & \texttt{data/calibration/} minimal content \\
\bottomrule
\end{longtable}


\subsection{What Already Exists}
\begin{itemize}[leftmargin=1.5em]
    \item Complete FPM acquisition software (controlled, swept, circular illumination patterns).
    \item Full FPM reconstruction pipeline including NaN-safe variant.
    \item RealESRGAN 4$\times$ super-resolution integrated.
    \item OV5647 sensor interface via \texttt{picamera2} over MIPI CSI-2.
    \item OLED illumination control via I2C (SSD1306, 128$\times$64).
    \item Real experimental data sets and JSON metadata schema.
    \item Calibration framework documented, though data remain sparse.
\end{itemize}

\subsection{What Is Missing}
\begin{itemize}[leftmargin=1.5em]
    \item Hardware design files (schematics, PCB, BOM).
    \item Power budget and conditioning design.
    \item Thermal analysis.
    \item Mechanical integration drawings.
    \item Communications/downlink software.
    \item Ground station software.
    \item Orbital or mission simulation artifacts.
    \item Radiation tolerance assessment.
    \item Watchdog/supervisor firmware.
    \item Flight software architecture and qualification plan.
\end{itemize}

\section{System Objective}

The INTISAT ISM payload shall demonstrate computational lensless microscopy from low Earth orbit by:
\begin{enumerate}[leftmargin=1.5em]
    \item acquiring raw intensity image sequences under controlled angular illumination patterns using an OLED or LED array at a fixed distance from a bare CMOS sensor,
    \item reconstructing a high-resolution complex-field image on-orbit using Fourier Ptychographic Microscopy,
    \item downlinking reconstructed science products and selected raw/diagnostic data,
    \item optionally applying super-resolution onboard or on the ground, and
    \item demonstrating the viability of computational imaging as a space instrument modality.
\end{enumerate}

\textbf{Primary objective:} demonstrate that a sub-100~g, sub-3~W computational microscope payload can produce sub-micron-equivalent imagery of biological, material, or optical samples in LEO, with onboard reconstruction and data compression.

\section{Mission Context and Assumptions}
\begin{longtable}{L{4cm} L{3cm} L{7cm}}
\toprule
\textbf{Parameter} & \textbf{Value} & \textbf{Justification / Source} \\
\midrule
Orbit type & LEO Sun-Synchronous & Stable thermal environment, regular ground contact \\
Altitude & 500--550 km & Typical CubeSat deployment altitude \\
Inclination & $\sim$97.5$^\circ$ & SSO at 500 km \\
Orbital period & $\sim$95 min & Derived from altitude \\
Ground passes/day & 4--6 useful & Typical LEO contact geometry \\
Mission duration & 12 months & Minimum demonstration mission \\
Platform & 3U CubeSat & 10$\times$10$\times$30 cm \\
Payload mass allocation & $\leq$300 g & Assumption \\
Payload power allocation & $\leq$5 W peak, $\leq$3 W avg & Assumption \\
Downlink band & UHF/VHF or S-band & TBD by bus architecture \\
Downlink rate & 9.6 kbps to 1 Mbps & Assumption; drives data budget \\
\bottomrule
\end{longtable}


\section{Payload Functional Description}

The payload operates according to the following chain:
\begin{enumerate}[leftmargin=1.5em]
    \item Command reception from OBC.
    \item Sample readiness check.
    \item Illumination sequencing over a 9$\times$9 angular grid.
    \item Raw image acquisition per angle with metadata logging.
    \item Onboard preprocessing: dark subtraction, flat-field correction, bad-pixel masking.
    \item FPM reconstruction into a high-resolution complex-field image.
    \item Optional super-resolution stage.
    \item Compression, packetization, storage, and downlink.
\end{enumerate}

One complete acquisition session consists of 81 raw frames, one reconstructed image, optional super-resolved output, and metadata.

\section{Subsystem Decomposition}

\subsection{Optical Subsystem}
\begin{longtable}{L{4cm}Y}
\toprule
\textbf{Field} & \textbf{Value} \\
\midrule
Purpose & Define lensless imaging geometry; sample placed directly above sensor \\
Inputs & Sample, illumination wavefront \\
Outputs & Diffracted intensity pattern on sensor plane \\
Design drivers & Working distance (8 mm), sample proximity, absence of objective lens \\
Dependencies & Mechanical integration, fixed z-distance \\
Risks & Contamination, stray light, z-drift due to thermal expansion \\
Still to design & Flight-compatible sample holder, dust protection, thermal compensation \\
\bottomrule
\end{longtable}


\subsection{Image Sensor Subsystem}
\begin{longtable}{L{4cm}Y}
\toprule
\textbf{Field} & \textbf{Value} \\
\midrule
Sensor & OmniVision OV5647, 1/4'' CMOS, 5.04 MP \\
Resolution & 2592$\times$1944 px, 1.4 $\mu$m pitch \\
Interface & MIPI CSI-2, 84 MHz pixel clock \\
Controlled via & \texttt{picamera2} with manual exposure/gain \\
Major risks & Radiation damage, thermal noise, hot-pixel accumulation \\
Current status & Bench-tested and operational \\
\bottomrule
\end{longtable}


\subsection{Illumination Subsystem}
\begin{longtable}{L{4cm}Y}
\toprule
\textbf{Field} & \textbf{Value} \\
\midrule
Current device & SSD1306 OLED (bench only) \\
Interface & I2C at 0x3C, 400 kHz \\
Emission & Approximately 530 nm \\
Pattern & 9$\times$9 grid, 81 positions \\
Major risk & OLED vacuum incompatibility and outgassing \\
Recommendation & Replace with discrete LED matrix for flight version \\
\bottomrule
\end{longtable}


\subsection{Compute Subsystem}
\begin{longtable}{L{4cm}Y}
\toprule
\textbf{Field} & \textbf{Value} \\
\midrule
Recommended compute core & Raspberry Pi Compute Module 4 \\
Processor & BCM2711, quad-core Cortex-A72 @ 1.5 GHz \\
RAM & 4 GB LPDDR4 \\
Storage & 32 GB eMMC \\
Software stack & Python 3.11, numpy, scipy, picamera2, fpm-py, optional PyTorch \\
FPM processing time & 5--8 min for 81 frames \\
Active power & $\sim$3.5--4.5 W \\
Key risks & SEU, thermal throttling, OS corruption \\
\bottomrule
\end{longtable}


\section{Hardware Architecture Options}
\begin{longtable}{L{2.2cm}L{4.2cm}L{6.3cm}}
\toprule
\textbf{Option} & \textbf{Description} & \textbf{Verdict} \\
\midrule
A & Full Raspberry Pi 4 board as-is & Rejected: excess connectors, poor CubeSat integration, weak novelty \\
B & CM4 + custom carrier board & Strong candidate: compact, tailored interfaces, better integration \\
C & CM4 + distributed support boards & Viable but more complex for first implementation \\
D & MCU supervisor + CM4 hybrid & Recommended: low-power standby, watchdog, stronger flight credibility \\
\bottomrule
\end{longtable}


\section{Recommended Architecture}

The recommended approach is \textbf{Option D}: a hybrid architecture with a low-power supervisor MCU and a CM4 compute core.

\begin{itemize}[leftmargin=1.5em]
    \item The MCU handles watchdog, housekeeping telemetry, power gating, and OBC interface.
    \item The CM4 powers on only during acquisition and processing windows.
    \item This reduces thermal load and average power consumption.
    \item It also makes the payload more defensible as a real CubeSat subsystem rather than a repackaged development board.
\end{itemize}

\section{Block Diagrams}

\subsection{Overall Payload Architecture}
\begin{lstlisting}[basicstyle=\ttfamily\small,frame=single]
SAMPLE INTERFACE
      |
      v
+-------------------+      +----------------------+
|    SENSOR BOARD   |<-----|  ILLUMINATION BOARD  |
| OV5647 (lensless) |      |  LED matrix + driver |
+---------+---------+      +----------+-----------+
          | CSI-2                     | I2C
          v                           v
+--------------------------------------------------+
|                 CM4 CARRIER BOARD                |
| +-------------------+  +-----------------------+ |
| | Raspberry Pi CM4  |  | STM32 Supervisor MCU  | |
| | BCM2711 + eMMC    |  | WDT, HK, Power Ctrl   | |
| +-------------------+  +-----------------------+ |
|        Power conditioning + current monitor      |
+---------------------------+----------------------+
                            |
                            v
                      SATELLITE OBC
\end{lstlisting}

\subsection{Data Flow}
\begin{lstlisting}[basicstyle=\ttfamily\small,frame=single]
LED pattern -> OV5647 capture -> TIFF + JSON metadata -> eMMC storage
       -> dark subtraction + flat field
       -> FPM reconstruction
       -> reconstructed image / phase image
       -> compression + packetization
       -> OBC interface -> RF downlink -> Ground station
\end{lstlisting}

\section{Imaging Chain Definition}
\begin{longtable}{L{4cm}L{4cm}L{5cm}}
\toprule
\textbf{Parameter} & \textbf{Value} & \textbf{Notes} \\
\midrule
Sensor & OV5647 & Raspberry Pi Camera Module v2 compatible \\
Full resolution & 2592$\times$1944 px & 5.04 MP \\
Pixel size & 1.4 $\mu$m & Physical pitch \\
Bit depth & 10-bit native, stored as 16-bit TIFF & Raw Bayer \\
Exposure range & 1 $\mu$s -- 6000 ms & Software controlled \\
Default exposure & 150 ms & From capture script \\
Working distance & 8.0 mm & Illumination-to-sensor \\
Illumination wavelength & 530 nm & Green dominant \\
FPM output resolution & 4$\times$ linear enhancement & Approx. 0.35 $\mu$m equivalent \\
Frames per session & 81 & 9$\times$9 grid \\
Acquisition time & 3--5 min & Includes settle delay and overhead \\
FPM reconstruction time & 5--8 min & CM4 CPU-bound estimate \\
\bottomrule
\end{longtable}


\section{Communications and Protocol Selection}
\begin{longtable}{L{3cm}L{3.5cm}L{3cm}L{2cm}L{4cm}}
\toprule
\textbf{Bus} & \textbf{Connection} & \textbf{Protocol} & \textbf{Rate} & \textbf{Justification} \\
\midrule
Sensor $\rightarrow$ CM4 & OV5647 to BCM2711 & MIPI CSI-2 & 84 MHz & Native high-speed camera interface \\
CM4 $\rightarrow$ LED driver & BCM2711 to LED driver & I2C & 400 kbps & Simple pattern control \\
CM4 $\leftrightarrow$ STM32 & Compute to supervisor & UART & 115200 baud & Sufficient for commands/HK \\
STM32 $\rightarrow$ OBC & Payload to spacecraft bus & CAN & 250 kbps & Fault-tolerant, standard CubeSat choice \\
\bottomrule
\end{longtable}


\section{Data Budget Tables}
\begin{longtable}{L{3.8cm}L{1.5cm}L{4cm}L{4cm}}
\toprule
\textbf{Item} & \textbf{Count} & \textbf{Raw Size} & \textbf{Compressed Size} \\
\midrule
Raw TIFF frame & 81 & 10.1 MB each ($\sim$818 MB/session) & $\sim$2 MB PNG each ($\sim$162 MB/session) \\
Dark frame & 1 & 10.1 MB & $\sim$2 MB \\
Metadata JSON & 81 + 1 & $\sim$344 KB total & --- \\
FPM reconstructed image & 1 & $\sim$45 MB TIFF & $\sim$2 MB JPEG2000 \\
Phase image & 1 & $\sim$45 MB & $\sim$5 MB \\
Preview thumbnail & 1 & $\sim$512 KB & $\sim$100 KB \\
\bottomrule
\end{longtable}


\section{Power Budget Tables}
\begin{longtable}{L{4cm}L{2cm}L{2cm}L{2cm}L{4cm}}
\toprule
\textbf{Component} & \textbf{Standby} & \textbf{Active} & \textbf{Peak} & \textbf{Notes} \\
\midrule
BCM2711 / CM4 & 0 W & 3.5 W & 4.5 W & Powered off in standby \\
OV5647 sensor & 0.01 W & 0.2 W & 0.3 W & Low-power CMOS \\
LED array & 0 W & 0.07 W & 0.12 W & One active LED example \\
STM32 supervisor & 0.01 W & 0.05 W & 0.1 W & Always available \\
Power conversion losses & 0.03 W & 0.8 W & 1.0 W & Buck/LDO losses \\
\midrule
\textbf{Total} & \textbf{$\sim$0.07 W} & \textbf{$\sim$5.3 W} & \textbf{$\sim$7.1 W} & \\
\bottomrule
\end{longtable}


\section{Calibration and Metrology Plan}
\begin{longtable}{L{3.8cm}L{3.8cm}L{2.5cm}L{2.5cm}L{3cm}}
\toprule
\textbf{Calibration Type} & \textbf{Method} & \textbf{Frequency} & \textbf{Storage} & \textbf{Notes} \\
\midrule
Dark frame & Same T, exposure, gain; LED off & Before each session + daily & eMMC & Must track sensor aging \\
Flat field & Uniform sample at each angle & Pre-launch, monthly & eMMC & Critical for quantitative imaging \\
PRNU & Uniform illumination, multiple exposures & Pre-launch & Lookup table & Fixed sensor characteristic \\
Thermal characterization & Dark frames vs temperature & Pre-launch TVAC & Lookup table & Temp-compensated correction \\
\bottomrule
\end{longtable}


\section{Validation and Verification Plan}
\begin{longtable}{L{1.8cm}L{3.5cm}L{4.2cm}L{4.2cm}L{2cm}}
\toprule
\textbf{Test ID} & \textbf{Objective} & \textbf{Method} & \textbf{Expected Output} & \textbf{Pass Criterion} \\
\midrule
OPT-01 & Verify lensless working distance & Known target at z = 8 mm & Sharp diffraction pattern & FOV within $\pm$0.5 mm \\
SEN-01 & Sensor verification & Functional camera test & Stable stream + control & All modes functional \\
CAL-01 & Dark frame calibration & Average 50 dark frames & Master dark file & Stable noise floor \\
CAL-02 & Flat-field calibration & Uniform illumination over all angles & 81 flat frames & Uniformity $>$95\% \\
IQ-01 & FPM quality & Reconstruct known target & Reconstructed image & Meets target resolution criterion \\
OBP-01 & End-to-end onboard chain & Capture to packetization test & Decodable science packet & Image quality acceptable \\
THM-01 & Thermal-vacuum operation & Full session at hot/cold extremes & Functional output & Completes successfully \\
VIB-01 & Launch vibration & Standard vibration test & Pre/post function check & No failures \\
\bottomrule
\end{longtable}


\section{Comparison Against Prior Work}

\begin{itemize}[leftmargin=1.5em]
    \item Many educational CubeSats fly cameras and downlink raw photos.
    \item PhiSat-class systems do onboard filtering/compression for Earth observation.
    \item Biology experiments in orbit usually rely on conventional microscopy or larger ISS-class platforms.
    \item Ground-based FPM exists, but the current differentiator is onboard FPM reconstruction in orbit for a lensless payload.
\end{itemize}

\section{Novelty Statement}

\emph{First demonstration of Fourier Ptychographic Microscopy executed entirely onboard a CubeSat in low Earth orbit, producing super-resolved complex-field microscopy images without an objective lens, using a custom carrier board based on the Raspberry Pi Compute Module 4.}

\section{Risks and Failure Points}
\begin{longtable}{L{1.5cm}L{6.2cm}L{2cm}L{2cm}L{4cm}}
\toprule
\textbf{ID} & \textbf{Description} & \textbf{Likelihood} & \textbf{Severity} & \textbf{Mitigation} \\
\midrule
R-01 & OLED vacuum incompatibility / outgassing & High & Critical & Replace OLED with LED matrix \\
R-02 & BCM2711 SEU during reconstruction & Medium & High & MCU watchdog, scrubbing, safe software \\
R-03 & OV5647 radiation degradation & High & Medium & In-flight dark-frame updates \\
R-04 & z-distance thermal drift & Medium & High & Low-CTE spacer + calibration \\
R-05 & eMMC corruption & Medium & High & Journaling, redundancy, safe writes \\
R-10 & No hardware design files exist yet & Certain & Critical & Start carrier PCB immediately \\
\bottomrule
\end{longtable}


\section{Next Engineering Actions}

\begin{itemize}[leftmargin=1.5em]
    \item Define OBC interface (CAN vs UART, packet format, command set).
    \item Select the flight illumination element (LED matrix strongly preferred).
    \item Define EPS interface and available current/voltage.
    \item Start CM4 carrier schematic in KiCad.
    \item Design sensor board and illumination board.
    \item Adapt acquisition code for headless flight operation.
    \item Implement telemetry/packetization and supervisor firmware.
\end{itemize}

\section*{Appendix A --- Draft Requirements}
\begin{longtable}{L{2.5cm}L{2.2cm}L{8cm}}
\toprule
\textbf{Req ID} & \textbf{Category} & \textbf{Requirement} \\
\midrule
PAY-REQ-001 & Functional & The payload shall acquire at least 49 raw intensity images per session under angularly diverse illumination. \\
PAY-REQ-002 & Functional & The payload shall reconstruct a high-resolution complex-field image onboard using FPM. \\
PAY-REQ-004 & Interface & The payload shall interface with the OBC via CAN bus at 250 kbps. \\
PAY-REQ-006 & Power & Average payload power shall not exceed 3 W at nominal duty cycle. \\
PAY-REQ-008 & Mass & Total payload mass shall not exceed 250 g. \\
PAY-REQ-009 & Volume & The payload shall fit within 1U (90$\times$90$\times$90 mm). \\
PAY-REQ-012 & Reliability & The payload shall include a hardware watchdog capable of resetting the CM4 if no heartbeat is received within 30 s. \\
PAY-REQ-015 & Calibration & The payload shall execute a dark-frame calibration sequence on command. \\
\bottomrule
\end{longtable}



\end{document}
